\subsubsection{Selection of Threshold $\vartheta_1$}
Threshold $\vartheta_1$ is used for pruning tree $\hDT_t$ in \cref{alg:ATC} to
obtain a data-adaptive tree $\hAT_t$ after invoking FOs for each layer.
We find that invoking FOs on the pruned tree may reduce estimation error if
threshold $\vartheta_1$ is carefully chosen.
To understand the reason, without loss of generality, let us consider a tree
without pruning and a tree pruned at node $a$, respectively (see
\cref{fig:branch_of_tree}).
If we spend the same amount of privacy budget $\epsilon_{t,2}$ and use the same
number of users to estimate the node frequencies in each layer of these two
trees, then the estimate of each node frequency will have the same variance
$\var(\epsilon_{t,2})$.
While for the pruned tree, properties of those pruned nodes are approximated by
recursively halving their parent nodes' properties.
Therefore, estimation errors are different only for the two subtrees rooted at
node $a$, excluding node $a$.
Assume the subtree has height $h+1$.

For the left tree, the total estimation error of the subtree rooted at node $a$
(excluding node $a$) is
\[
  \err_1 \triangleq
  \sum_{i=1}^h 2^i\E[(\hat{f_i} - f_i)^2]
  = 2(2^h-1)\var(\epsilon_{t,2}).
\]

For the right tree, the total estimation error of the subtree rooted at node $a$
(excluding node $a$) is
\begin{align*}
  \err_2
  &\triangleq \sum_{i=1}^h 2^i\E[(\hat{f}_i-f_i)^2]
    = \sum_{i=1}^h 2^i\E[(\frac{\hat{f}_0}{2^i}-f_i)^2] \\
  &= \sum_{i=1}^h 2^i\E( \frac{\hat{f}_0^2}{2^{2i}} - \frac{\hat{f}_0f_i}{2^{i-1}}
    +f_i^2 )
  = \sum_{i=1}^h 2^i(\frac{\var(\epsilon_{t,2}) + f_0^2}{2^{2i}}
    - \frac{f_0f_i}{2^{i-1}} + f_i^2) \\
  &\leq \sum_{i=1}^h 2^i(\frac{\var(\epsilon_{t,2})+\vartheta_1^2}{2^{2i}}+\vartheta_1^2)
    = \sum_{i=1}^h(\frac{\var(\epsilon_{t,2})+\vartheta_1^2}{2^i} + 2^i\vartheta_1^2) \\
  &\leq \var(\epsilon_{t,2}) + \vartheta_1^2 + 2(2^h-1)\vartheta_1^2 \\
  &= \var(\epsilon_{t,2}) + (2^{h+1}-1)\vartheta_1^2.
\end{align*}
If we require the estimation error for the right tree is no larger than the left
tree, i.e., $\err_2\le \err_1$, then we obtain
\begin{equation}\label{eq:theta1}
  \vartheta_1
  \leq \sqrt{\big(\frac{2^{h+1}-3}{2^{h+1}-1}\big)\cdot\var(\epsilon_{t,2})}
  \leq \sqrt{\var(\epsilon_{t,2})}.
\end{equation}
Therefore, if threshold $\vartheta_1$ is chosen according to
Condition~\eqref{eq:theta1}, the pruned tree will have smaller overall
estimation error than the full binary tree, thus demonstrating the usefulness of
tree pruning operation in \cref{alg:ATC}.
\begin{figure}[htp]
  \centering
  \begin{tikzpicture}[
  thick, font=\small, inner sep=0pt,
  nd/.style={circle,minimum size=2.5mm},
  nd0/.style={nd,draw,densely dotted},
  nd1/.style={nd,draw},
  line0/.style={densely dotted},
  line1/.style={},
  ]

  %% left tree
  \coordinate (O1);
  \foreach \i in {0,...,7}{
    \node[nd1,right=.4*\i of O1] (a3\i) {};
  }

  \foreach \i in {0,...,3}{
    \path (a3\the\numexpr\i*2\relax) -- (a3\the\numexpr\i*2+1\relax)
    node[midway,above=.3,nd1] (a2\i) {};

    \draw (a3\the\numexpr\i*2\relax) -- (a2\i) -- (a3\the\numexpr\i*2+1\relax);
  }

  \foreach \i in {0,1}{
    \path (a2\the\numexpr\i*2\relax) -- (a2\the\numexpr\i*2+1\relax)
    node[midway,above=.3,nd1] (a1\i) {};

    \draw (a2\the\numexpr\i*2\relax) -- (a1\i) -- (a2\the\numexpr\i*2+1\relax);
  }
  \path (a10) -- (a11) node[midway,above=.3,nd1] (a0) {};
  \draw (a10) -- (a0) -- (a11);

  % layer number
  \node at (a11) {$a$};
  \node[right=.7 of a11] (l0) {$0$};
  \node at (l0 |- a20) (l1) {1};
  \node at (l0 |- a30) (l2) {$2$};
  \node[below=0 of l2] (ld) {$\vdots$};
  \node[below=.1 of ld] (lh) {$h$};
  \node at (a11 |- ld) {$\vdots$};


  %% right tree
  \coordinate[right=4 of O1] (O2);
  \foreach \i/\j in {0/1,1/1,2/1,3/1,4/0,5/0,6/0,7/0}{
    \node[nd\j,right=.4*\i of O2] (b3\i) {};
  }

  \foreach \i/\j in {0/1,1/1,2/0,3/0}{
    \path (b3\the\numexpr\i*2\relax) -- (b3\the\numexpr\i*2+1\relax)
    node[midway,above=.3,nd\j] (b2\i) {};

    \draw[line\j] (b3\the\numexpr\i*2\relax) -- (b2\i) --
    (b3\the\numexpr\i*2+1\relax);
  }

  \foreach \i/\j in {0/1,1/0}{
    \path (b2\the\numexpr\i*2\relax) -- (b2\the\numexpr\i*2+1\relax)
    node[midway,above=.3,nd1] (b1\i) {};

    \draw[line\j] (b2\the\numexpr\i*2\relax) -- (b1\i) --
    (b2\the\numexpr\i*2+1\relax);
  }

  \path (b10) -- (b11) node[midway,above=.3,nd1] (b0) {};
  \draw (b10) -- (b0) -- (b11);

  % \path (b0) -- (b11) coordinate[pos=.3] (A) coordinate[pos=.7] (B);
  % \draw (b0) -- (A);
  % \draw (b11) -- (B);
  % \draw[densely dotted] (A) -- (B);

  % layer number
  \node at (b11) {$a$};
  \node[right=.7 of b11] (l0) {$0$};
  \node at (l0 |- b20) (l1) {1};
  \node at (l0 |- a30) (l2) {$2$};
  \node[below=0 of l2] (ld) {$\vdots$};
  \node[below=.1 of ld] (lh) {$h$};
  \node at (b11 |- ld) {$\vdots$};

\end{tikzpicture}

  \caption{Example of a same tree without pruning (left) and pruning at node $a$
  (right)}
  \label{fig:branch_of_tree}
  \Description{This figure illustrates a comparison between a tree structure before and after pruning. The tree on the left shows the original structure, where all branches are intact.}
\end{figure}